
Five research questions are addressed, corresponding to five validated sub-hypotheses:

\begin{itemize}
\item \textbf{RQ1 (H-M1)}: Do benchmark trajectories exhibit non-random decelerating structure?
\item \textbf{RQ2 (H-E1, H-M2)}: Is the logistic model statistically preferred over alternatives?
\item \textbf{RQ3 (H-M3)}: Are fitted parameters physically plausible?
\item \textbf{RQ4 (H-M4)}: Does the dual-criterion detector recover community-recognized saturation dates within $\pm 6$ months?
\item \textbf{RQ5 (H-C1)}: Does the pipeline generalize to 30--49-entry benchmarks?
\end{itemize}

\subsubsection{4.1 Benchmarks}

\begin{table*}[htbp]
\centering
\resizebox{\dimexpr 2\columnwidth+\columnsep\relax}{!}{%
\begin{tabular}{lllll}
\toprule
Benchmark & Release & Post-processing Entries & Ground Truth Saturation Date & Ground Truth Source \\
\midrule
GLUE & Feb 2019 & 51 (deduplicated) & Sept 2019 & SuperGLUE paper (Wang et al., 2019b) \\
SuperGLUE & May 2019 & 50 (deduplicated) & Jun 2021 & BIG-bench (Srivastava et al., 2022) \\
\bottomrule
\end{tabular}
}
\end{table*}

\textbf{Data source}: Curated historical scores compiled from published model papers (BERT, XLNet, RoBERTa, ALBERT, MT-DNN, T5, DeBERTa, etc.). The Papers With Code API (paperswithcode-client v0.3.1) is non-functional as of the time of this work --- all endpoints redirect to HuggingFace (HTTP 302). The curated data faithfully represents published score progressions documented in the original benchmark papers and major model papers, and produces S-curves consistent with the literature. The implications of this data infrastructure limitation are discussed in Section 6.2.

\subsubsection{4.2 Baselines for RQ2}

\textbf{Linear} ($y = at + b$): null hypothesis of constant-rate improvement. \textbf{Power-law} ($y = at^b$): sub-linear alternative with no asymptote. Both are standard growth models; inclusion makes the AIC comparison a genuine three-way test.

\subsubsection{4.3 Synthetic Benchmarks for RQ5}

20 synthetic timeseries at n $\in$ {30, 35, 40, 45} entries (5 per count), logistic parameters drawn from GLUE/SuperGLUE ranges, Gaussian noise ($\sigma$ = 0.02), random seed 42. Synthetic data was used because the Papers With Code \texttt{benchmark\textbackslash \_list} API method was unavailable.

\subsubsection{4.4 Evaluation Metrics}

$R^2$ (fit quality), $\Delta AIC$ (model preference), Spearman $\rho$ (temporal structure), saturation date error in months (application accuracy), convergence rate and mean $R^2$ at low N (generalization).

\subsubsection{4.5 Ground Truth Validity Note}

The saturation dates used for validation are community-documented events recorded retroactively in published papers, not independently measured saturation timestamps. This creates an inherent circularity: a detector designed to systematize community judgment is validated against that same community judgment. An ideal validation would employ blind annotation or prospective prediction prior to community action. This structural limitation is reported in Section 6.2 (L6).

